% ===========================================================================
\section{Smart Contract Design and Implementation}\label{sec:contracts}
% Source material: contracts/src and contracts/test, which match
% the deployed sources.
% ===========================================================================

The on-chain part of \projectname{} is six Solidity files, 632 lines with
their comments, compiled with Solidity 0.8.28~\cite{solidity}. They build on
OpenZeppelin Contracts 5.6~\cite{openzeppelin} for access control, clones and
Merkle proofs, and on Semaphore's verifier, version 4.14. \texttt{Election} is an abstract contract that inherits OpenZeppelin's
\texttt{AccessControl} and \texttt{Initializable} and our
\texttt{VoterRegistry}. \texttt{MerkleElection} and \texttt{SemaphoreElection}
extend it with a \texttt{castVote} each.

\subsection{Contract responsibilities}

\subsubsection*{\texttt{ElectionFactory}}

The only way to create an election. Its state is the address of the
implementation it clones, fixed at deployment, and a mapping
\texttt{isElection}.
\texttt{createElection} can be called by anyone. It deploys a clone, records
it, and initializes it in the same transaction, so nobody can initialize a
clone before its creator does; the caller becomes the election's authority.
The implementation itself disables its initializer in its constructor, so it
can never be configured and used as an election. \texttt{isElection} tells a
genuine election from a look-alike, and the \texttt{ElectionCreated} events
list every election the factory made, so the factory stores no list of its
own. We deploy two factories, one per tier.

\subsubsection*{\texttt{Election} and \texttt{VoterRegistry}}

\texttt{Election} holds one election's lifecycle, candidates, ballots and
tally, and each clone has its own copy of this state:

\begin{itemize}
  \item the phase, the schedule (four timestamps packed into one storage
        slot), the title and the IPFS address of the ballot description;
  \item the candidates (1 to 64, each with a name and an optional manifesto
        address), the counting rule, the tally and, once finalized, the
        winners;
  \item \texttt{ballotsCast} and a mapping of used ballot keys: the voter's
        address in tier 2, the nullifier in tier 3.
\end{itemize}

\texttt{VoterRegistry}, an abstract base of \texttt{Election}, adds the voter
list: its Merkle root, the number of voters and its IPFS address. The three
are written together, once, and never change. No voter data is stored
on-chain.

Two roles are granted when the election is initialized:
\texttt{AUTHORITY\_ROLE} to its creator and \texttt{REGISTRAR\_ROLE} to the
registrar it names. Nobody holds the admin role, so neither role can be
granted to anyone else or revoked; a holder can only renounce its own.
Table~\ref{tab:entrypoints} lists every function that changes state, who may
call it and when.

\begin{table}[htbp]
  \centering
  \caption{State-changing functions of an election, in lifecycle order. Each
    time is one of the four in the election's schedule, compared with the
    block's timestamp.}
  \label{tab:entrypoints}
  \small
  \begin{tabular}{@{}L{3.0cm}L{2.3cm}L{6.4cm}L{5.4cm}@{}}
    \toprule
    \textbf{Function} & \textbf{Caller} & \textbf{Allowed when} & \textbf{Effect} \\
    \midrule
    \texttt{configure} & Authority & Created, before \mbox{\emph{registrationStart}} & Replaces the title, description and schedule \\
    \addlinespace[2pt]
    \texttt{setCandidates} & Authority & Created or Registration, before \mbox{\emph{votingStart}} & Replaces the candidates \\
    \addlinespace[2pt]
    \texttt{cancel} & Authority & Created or Registration, except when the list is published and the voting window is open & Moves to Cancelled \\
    \addlinespace[2pt]
    \texttt{openRegistration} & Anyone & Created, from \mbox{\emph{registrationStart}} & Moves to Registration \\
    \addlinespace[2pt]
    \texttt{publishVoterList} & Registrar & Registration, from \mbox{\emph{registrationEnd}} until \mbox{\emph{votingStart}}; once; no empty value & Writes the root, the voter count and the list's address \\
    \addlinespace[2pt]
    \texttt{openVoting} & Anyone & Registration, from \mbox{\emph{votingStart}} until \mbox{\emph{votingEnd}}, with the list published & Sets up the tally; moves to Voting \\
    \addlinespace[2pt]
    \texttt{castVote} & Anyone (tier 3); a listed wallet (tier 2) & Voting, before \mbox{\emph{votingEnd}} & Records one ballot \\
    \addlinespace[2pt]
    \texttt{closeVoting} & Anyone & Voting, from \mbox{\emph{votingEnd}} & Moves to Tallying \\
    \addlinespace[2pt]
    \texttt{finalize} & Anyone & Tallying & Records the winners; moves to Finalized \\
    \bottomrule
  \end{tabular}
\end{table}

A schedule is accepted only if its four times are in order, and a title only
if it is not empty. The candidates can change until voting is due, so the
meaning of a ballot never changes after it is cast. Both the schedule and the
candidates are limited by the clock as well as by the phase, because
registration and voting are open from their start times even before anyone
has recorded that on-chain.

\subsubsection*{\texttt{MerkleElection} (tier 2)}

Adds \texttt{castVote(choice, proof)}. The proof shows that the sender's
address is a leaf under the published root, using OpenZeppelin's
\texttt{MerkleProof}. Leaves are hashed exactly as OpenZeppelin's
\texttt{StandardMerkleTree} hashes an address, so anyone can rebuild the tree
from the published list with that library. The ballot key is the sender's
address, so the chain shows which wallet chose which candidate, and the voter
pays the gas.

\subsubsection*{\texttt{SemaphoreElection} (tier 3)}

Adds \texttt{castVote(choice, nullifier, points)}, where \texttt{points} is a
Groth16 proof~\cite{groth16}, and the read functions \texttt{scope} and
\texttt{merkleTreeDepth}. Its only state of its own is the verifier's address,
fixed at deployment, which clones read from the implementation's code.
\texttt{castVote} never reads the sender, so it does not matter who submits a
ballot.
Section~\ref{sec:zk} explains the scheme and Listing~\ref{lst:castvote} shows
the code.

\subsubsection*{\texttt{TallyLib}}

An internal library, compiled into each election contract, with the two
counting rules. Under plurality a choice is a candidate's index. Under approval
it is a bitmask with one bit per approved candidate: at least one bit must be
set, none beyond the last candidate, and each approved candidate gains one
vote. Under both rules the winners are every candidate tied for the most
votes, in ascending order, and there are none if nobody voted. The rule is set
when the election is created and has no setter.

\subsubsection*{\texttt{SemaphoreVerifier}}

Semaphore's Groth16 verifier, deployed unmodified. It holds one verification key for each
tree depth from 1 to 32 and has no state or owner.

\subsection{Election lifecycle}\label{sec:lifecycle}

Figure~\ref{fig:lifecycle-final} shows the lifecycle exactly as the modifiers
in \texttt{Election.sol} enforce it. Its steps fall into four kinds.

\begin{figure}[tb]
  \centering
  \resizebox{0.92\textwidth}{!}{% DVote election lifecycle as implemented in contracts/src/Election.sol.
% Bare tikzpicture: \input this inside a figure environment.
% Redrawn from figures/lifecycle.tex (kept unchanged for the Evaluation 1
% sources) with the exact guards of each function: who may call it, and
% between which of the four scheduled times.
\begin{tikzpicture}[
    font=\small,
    nohyph/.style   = {execute at begin node={\hyphenpenalty=10000\exhyphenpenalty=10000}},
    phase/.style    = {draw=black, thick, rounded corners=3pt, fill=white,
                       align=center, inner sep=4pt, text width=5.8cm,
                       minimum height=1.0cm, nohyph},
    terminal/.style = {phase, double, double distance=1.1pt},
    aborted/.style  = {draw=black, thick, rounded corners=3pt, dashed,
                       fill=black!7, align=center, inner sep=4pt,
                       text width=2.4cm, minimum height=1.0cm},
    trans/.style    = {-{Latex[length=2.4mm,width=2mm]}, draw=black,
                       line width=1.1pt},
    abort/.style    = {-{Latex[length=2.2mm,width=1.8mm]}, draw=black!65,
                       semithick, dash pattern=on 2.5pt off 2pt,
                       rounded corners=3pt},
    tlab/.style     = {align=left, inner sep=2pt, text width=6.6cm,
                       font=\footnotesize, nohyph},
  ]

  % ---- phases ----
  \node[phase]    (created)  at (6.0, 0)
    {\textbf{Created}\\the authority may change the schedule until \emph{registrationStart}};
  \node[phase]    (regn)     at (6.0, -2.5)
    {\textbf{Registration}\\the registrar publishes the list once,\\from \emph{registrationEnd} until \emph{votingStart}};
  \node[phase]    (voting)   at (6.0, -5.0)
    {\textbf{Voting}\\ballots accepted until \emph{votingEnd}};
  \node[phase]    (tallying) at (6.0, -7.5)
    {\textbf{Tallying}\\no more ballots};
  \node[terminal] (final)    at (6.0, -10.0)
    {\textbf{Finalized}\\tally and winners permanent};
  \node[aborted]  (cancelled) at (0.3, -1.25) {\textbf{Cancelled}\\no result};

  % ---- forward steps ----
  \draw[trans] (created.south)  -- (regn.north);
  \draw[trans] (regn.south)     -- (voting.north);
  \draw[trans] (voting.south)   -- (tallying.north);
  \draw[trans] (tallying.south) -- (final.north);

  \node[tlab, anchor=west] at (9.5, -1.25)
    {\texttt{openRegistration}: anyone, from \emph{registrationStart}.};
  \node[tlab, anchor=west] at (9.5, -3.75)
    {\texttt{openVoting}: anyone, from \emph{votingStart} until \emph{votingEnd},
     and only once the voter list is published.};
  \node[tlab, anchor=west] at (9.5, -6.25)
    {\texttt{closeVoting}: anyone, from \emph{votingEnd}. Ballots stop at
     \emph{votingEnd} whether or not anyone has called it.};
  \node[tlab, anchor=west] at (9.5, -8.75)
    {\texttt{finalize}: anyone, at any time. Records the winners.};

  % ---- cancellation ----
  \draw[abort] (created.west) -- (2.2, 0) -- (2.2, -0.85) -- (cancelled.east |- 0, -0.85);
  \draw[abort] (regn.west)    -- (2.2, -2.5) -- (2.2, -1.65) -- (cancelled.east |- 0, -1.65);
  \node[align=center, anchor=north, text width=3.4cm, inner sep=2pt,
        font=\footnotesize, nohyph] at (0.3, -1.95)
    {\texttt{cancel}: the authority only, from Created or Registration,
     unless the list is published and the voting window has begun};

  % ---- legend ----
  \draw[trans] (9.6, -9.8) -- (10.6, -9.8);
  \node[anchor=west, font=\footnotesize] at (10.7, -9.8) {anyone may call it once its time comes};
  \draw[abort] (9.6, -10.3) -- (10.6, -10.3);
  \node[anchor=west, font=\footnotesize] at (10.7, -10.3) {the authority only};
\end{tikzpicture}
}
  \caption{Election lifecycle as implemented. Each forward step names who may
    call it and from when; each phase names what happens inside it.}
  \label{fig:lifecycle-final}
\end{figure}

\begin{itemize}
  \item \emph{Time-gated and explicit.} Opening registration, opening voting
        and closing voting each need their scheduled time to have passed and a
        transaction from someone. Anyone may send it. Opening voting also has a
        deadline: once \emph{votingEnd} has passed, voting can never open.
  \item \emph{Explicit only.} \texttt{finalize} has no time of its own. It is
        possible as soon as voting has been closed, which itself needs
        \emph{votingEnd} to have passed. \texttt{cancel} waits for no scheduled
        time either: it needs the authority and the right phase, and is refused
        only while a published list's voting window is open.
  \item \emph{By the clock alone.} Some things happen without any
        transaction. Ballots stop being accepted at \emph{votingEnd} even while
        the phase still reads Voting, because \texttt{castVote} checks the time
        itself. The authority loses the right to change the schedule at
        \emph{registrationStart} and the candidates at \emph{votingStart},
        whether or not the phase has moved.
  \item \emph{Irreversible.} No function moves an election to an earlier
        phase, and the tests check this. The voter list can be published only
        once, and is frozen while the election is still in Registration.
        Finalized and Cancelled are final: no function that changes state
        accepts an election in either phase.
\end{itemize}

Everything after configuration is open to anyone, so an authority that
dislikes how the vote is going gains nothing by refusing to act. What the
authority keeps is the ability to cancel, and only while that cannot stop a
vote already under way. An election whose registrar never published a list,
or whose voting window passed without anyone opening voting, can still be
cancelled, so it does not stay open for good. In practice, if registration is
not yet open on-chain when the registrar publishes, the registrar opens it
first, and the election page offers the next step as a button to anyone with a
wallet.

\subsection{Anonymous eligibility and nullifiers}\label{sec:zk}

Tier 3 uses Semaphore version~4~\cite{semaphore}, with our own tree in place of
Semaphore's on-chain groups.

\paragraph{Identity.} The booth creates a random private key. From it
Semaphore derives a secret scalar $s$ and a public key $(A_x, A_y)$ on the
Baby Jubjub curve, and the identity's commitment is the Poseidon hash of that
public key~\cite{semaphorecircuit}. The commitment is the only value that
leaves the device. The secret stays in browser storage and in the backup file.

\paragraph{The voter list.} When registration closes, the registrar sorts the
commitments in ascending order, removes repeats and builds a Lean Incremental
Merkle Tree with Poseidon~\cite{leanimt}. In this binary tree a node with no
right child takes its left child's value, so $n$ leaves give a tree of depth
$\lceil \log_2 n \rceil$ with no padding. Sorting means the tree does not
depend on the order in which people registered, and a commitment's position
says nothing about when its owner registered. The list goes to IPFS and only
the root, the count and the list's address go on-chain.

\paragraph{The proof.} To vote, the booth rebuilds the tree from the published
list and proves, with Semaphore's circuit, that it knows a secret $s$ whose
commitment is a leaf under the root. The proof has four public values, and the
contract builds all four itself rather than taking them from the caller:
\[
  \bigl[\,\mathit{voterRoot},\ \mathit{nullifier},\ H(\mathit{choice}),\
  H(\mathit{scope})\,\bigr],
  \qquad H(x) = \mathrm{keccak256}(x) \gg 8,
\]
where the shift keeps a hash inside the proof system's field, as Semaphore
does. The scope is the election's own address. The contract also chooses the
verification key, through \texttt{merkleTreeDepth}: the depth is
$\lceil \log_2 n \rceil$ for the published count $n$, but at least 1, the
smallest depth Semaphore's circuits support. The caller supplies only the
choice, the nullifier and the proof, and a proof made for any other choice,
election, list or depth fails.

\paragraph{The nullifier.} The circuit computes the nullifier as
$\mathrm{Poseidon}(H(\mathit{scope}), s)$~\cite{semaphorecircuit}, and the
proof shows it was computed correctly. It is the same every time the same
secret votes in the same election, so \texttt{\_recordBallot} rejects a second
ballot as \texttt{AlreadyVoted}. It differs between elections, because the
scope is the election's address, so ballots from one voter in two elections
cannot be linked. And it cannot be turned back into the secret. A relayer
cannot change the choice, which is bound into the proof, or replay a ballot,
whose nullifier is already used.

\paragraph{The ballot cap.} Independently of the proof, \texttt{\_recordBallot}
refuses any ballot once \texttt{ballotsCast} has reached the published voter
count. This check does not depend on the trusted setup, so it still bounds the
number of ballots if the setup were ever broken.

What the chain records for each ballot is the \texttt{VoteCast} event with the
nullifier and the choice, plus the transaction's sender and time. The sender is
the relayer, or whichever wallet submitted the ballot. Nothing in it names the
voter; how much that hides depends on the size of the list and the timing,
as Section~\ref{sec:trust} explains.

\subsection{Gas-efficiency measures}\label{sec:gas-measures}

Section~6 measures the ones that can be measured. Unless marked as Sepolia's, every figure quoted here is under mainnet's current rules, from the
gas snapshot committed in \texttt{contracts/snapshots/} or the sweep in
\texttt{bench/results/gas.csv}; Table~\ref{tab:a1} lists the main ones.

\begin{itemize}
  \item \emph{The voter list stays off-chain.} Publishing a list writes a root,
        a count and an IPFS address in one transaction, and emits no list. It
        costs 147,165 gas in either tier at every list size
        from 2 to 65,536 voters, which parameter C1 confirms
        (Figure~\ref{fig:c1}). Registering costs voters no gas at all.
  \item \emph{Clones.} Creating a wallet election, a clone set up with three
        candidates, costs 595,090 gas. Deploying a full copy of the election
        contract instead costs 2,080,541 before it is even set up, and under
        Glamsterdam the same deployment cost 14,630,061
        (Table~\ref{tab:deployment}).
  \item \emph{Descriptions and manifestos on IPFS.} The chain holds only their
        IPFS addresses, with the title and the candidates' names, which a reader
        needs to make sense of the tally. Section~6 does not measure this
        separately.
  \item \emph{A constant-time tree depth.} The depth comes from one storage
        read and \texttt{Math.log2}, rather than a loop over the levels, which
        would add about 210 gas per level. A tier 3 ballot costs between
        296,634 and 296,670 gas at every list size from 4 to 65,536 voters
        (296,269 at 2), a spread of 36 gas.
  \item \emph{Small ballots.} An approval ballot is one 256-bit number,
        however many candidates it approves. The four schedule times share one
        storage slot. A ballot writes one new storage slot, for its key, and
        updates the ballot count and the tally.
  \item \emph{No forwarder.} A relayed ballot is a plain call to
        \texttt{castVote}. Forwarding it with a signature from a fresh key
        cost 127,194 gas more per ballot on Sepolia
        (Section~\ref{sec:req-changes}).
\end{itemize}

Some costs are kept on purpose. Titles and candidate names are stored as strings so
that the contract alone explains its tally, and roles go through
\texttt{AccessControl} rather than fixed addresses, because a clone cannot have
its own immutable values. The first ballot also costs more than later ones
(114,397 against 80,221 gas in tier 2, 330,870 against 296,670 in tier 3),
because it moves \texttt{ballotsCast} and a candidate's count away from zero.

\subsection{Selected implementation}

Listing~\ref{lst:castvote} is the tier 3 ballot function and the two values it
derives. Listing~\ref{lst:recordballot} is the part both tiers share. Both are
taken directly from the deployed source.

\lstinputlisting[language=Solidity, firstline=34, lastline=72, firstnumber=34,
  columns=fullflexible, keepspaces=true, float=tbp,
  caption={\texttt{SemaphoreElection.castVote} and the values it derives
    (\texttt{contracts/src/SemaphoreElection.sol}).},
  label={lst:castvote}]{../contracts/src/SemaphoreElection.sol}

\lstinputlisting[language=Solidity, firstline=262, lastline=284, firstnumber=262,
  columns=fullflexible, keepspaces=true, float=tbp,
  caption={How every ballot is accepted and counted, in either tier
    (\texttt{contracts/src/Election.sol}).},
  label={lst:recordballot}]{../contracts/src/Election.sol}

The code shows several properties directly. The public values of the proof
are built from storage and from the contract's own address; the only inputs
taken from the caller are the choice, the nullifier and the proof. The one
external call, to the verifier, is to a \texttt{view} function, so it runs as a
static call and cannot change any state. All state changes come after it.
\texttt{\_recordBallot} checks the key, the choice and the cap before writing
anything, so a ballot is either fully counted or has no effect. The same
function serves both tiers, so the double-vote and cap checks cannot drift
apart between them.

\subsection{Testing and static analysis}\label{sec:contract-tests}

The contracts have 90 Foundry~\cite{foundry} tests in six suites
(Table~\ref{tab:contract-tests}), and all pass at the deployed commit. Five are
fuzz tests, which Foundry runs with 256 random inputs each by default and
1,000 in continuous integration. There are no stateful invariant tests.

\begin{table}[htbp]
  \centering
  \caption{Contract test suites.}
  \label{tab:contract-tests}
  \small
  \begin{tabular}{@{}L{3.3cm}R{1.2cm}L{13.0cm}@{}}
    \toprule
    \textbf{Suite} & \textbf{Tests} & \textbf{What they assert} \\
    \midrule
    \texttt{Election} & 47 & Every function's caller and time window from Table~\ref{tab:entrypoints}; phases never go back; every cancellation case; the list is published once and then frozen; a key votes once (fuzzed over arbitrary keys); ballots stop at the voter count; approval ballots, ties and an election nobody voted in; roles cannot be granted or revoked by anyone \\
    \addlinespace[2pt]
    \texttt{SemaphoreElection} & 17 & Real Semaphore proofs are accepted and counted; a second ballot, a changed choice, a ballot for another election, a proof against another root, a tampered proof and a proof at the wrong depth are all refused; the event carries only the nullifier and the choice; any account may submit a ballot, and none can replay a counted one; the depth is the smallest that fits the list (fuzzed over voter counts up to $2^{32}-1$) \\
    \addlinespace[2pt]
    \texttt{MerkleElection} & 10 & A listed wallet is counted once; a wallet not on the list fails with any proof (fuzzed over addresses), as does someone else's proof; leaves match OpenZeppelin's tree, and proofs built in TypeScript are accepted \\
    \addlinespace[2pt]
    \texttt{ElectionFactory} & 7 & The caller becomes the authority; every election is recorded; clones are independent and cannot be initialized twice; zero addresses are refused; the creation event names the authority and the registrar \\
    \addlinespace[2pt]
    \texttt{TallyLib} & 7 & Winners have the most votes and every tied leader is returned, in order (fuzzed); an approval ballot adds exactly one vote per approved candidate (fuzzed); what counts as a valid choice under each rule \\
    \addlinespace[2pt]
    \texttt{GasSnapshot} & 2 & Runs a full election in each tier and records the gas of every step; continuous integration fails if any figure moves from the committed snapshot \\
    \bottomrule
  \end{tabular}
\end{table}

The tier 3 tests do not use mock proofs. The ballots in
\texttt{test/fixtures/semaphore-ballots.json} are real Semaphore proofs, made by
the same TypeScript code the booth runs, so the browser's proofs and the
on-chain verifier are tested against each other. The tier 2 tests likewise use
a voter list built by the TypeScript tree builder.

\texttt{forge coverage} reports every line (162 of 162), statement (174 of
174), branch (32 of 32) and function (43 of 43) of the contracts covered.

Static analysis runs on every pull request and every push to the main
branch. Continuous integration checks formatting,
runs \texttt{forge lint} with warnings treated as errors, and runs
Slither~\cite{slither} 0.11.6, which fails the build on any finding. Slither
analyses our six files only: OpenZeppelin and Semaphore's verifier are
excluded as dependencies. On the deployed sources it reports no findings, with
these dispositions for what is suppressed:

\begin{itemize}
  \item \emph{Block timestamp} (Slither's \texttt{timestamp} detector, rated
        low impact; excluded in \texttt{slither.config.json}). It flags every
        comparison with \texttt{block.timestamp}, and the lifecycle consists of
        such comparisons by design. On proof-of-stake Ethereum a block's
        timestamp must equal the start time of its slot~\cite{bellatrix}, so a
        block proposer cannot choose it. The most a proposer can do is skip its
        slot, which delays the next block by one slot. Against windows of
        minutes to days, that cannot move a phase meaningfully. We judge every
        instance a false positive.
  \item \emph{Lints.} \texttt{forge lint}'s block-timestamp lint is excluded
        for the same reason as Slither's detector. Four lines suppress its
        reentrancy-events lint, each with a comment saying why it does not
        apply: before each of those events the only call is to create a
        clone, which runs no code, to the read-only verifier, or none at all.
\end{itemize}

What this does not give us is an independent audit or formal verification.
The tests show that the contracts do what we intended; they cannot show that
we intended the right things.
